% Pillar-5 (MIN-REPORT-RL) limitations
\section{Limitations}
\label{sec:p5-limitations}
Position papers owe their readers an honest account of the evidence's edges;
five limitations are load-bearing here. First, the headline exhibits
(the $17\times$ under-specified same-label comparison, the DAPO-label exhibit, the
12-cell head-to-head, and the 368-run audit) are internal program records:
they were run under a declared protocol and are summarized faithfully, but
they have not undergone the full artifact-release pipeline of the \ours{}
Stratum-A results, and the closed stack involved cannot be independently
re-executed by readers---an irony we acknowledge and which the manifest
standard is designed to at least make legible. Second, item 4's theory has no
gradient-level support: the earlier toy correlation ($+0.71$) cannot be
recomputed, rests on about 30 serially dependent rows from one run, and is
largely mechanical, so we no longer cite it. Third, our evidence is concentrated on GSM8K-style
verifiable-reward tasks and the GRPO family; the eight-item standard (seven
run-manifest fields plus pass@$k$ reporting) was chosen for that regime, and
reward-model-based RLHF stacks may need additional items
(reward-model version and training data at minimum). Fourth, the flip-risk
grades encode an ordering (outcome-critical fields such as the base
checkpoint above algorithm fields, above hyperparameters and runtime) that is
motivated by our own exhibits but not validated against a broad external
corpus; the field-to-grade assignments, especially R2 vs.\ R3, will need
community calibration. Fifth, the toolchain of \secref{sec:p5-toolchain} is a
specification, not a released product suite; compliance overhead has not
been measured and manifest adapters remain to be built. The retained
exhibits motivate disclosure, not a universal ranking of stack effects
above algorithm effects.
